% LaTeX companion of hw10.typ. Both formats are editable.
\section{Homework 10 --- submitted work}

\chapter{Part A}

The assigned exercises are 6.1: 20, 54; 6.2: 42, 50; 6.3: 14; and 7.1: 12, 18, 42.

\section{6.1 Exercise 20}

By the row operations shown in the submission, the determinant of the given \(k \times k\) matrix reduces to

\(\det(A) = 1\)

for every \(k\).

\section{6.1 Exercise 54}

The answer uses positivity to conclude that the determinant of the displayed \(5 \times 5\) matrix is positive, and records

\(\det(A) = 1000.\)

\section{6.2 Exercise 42}

For a QR factorization \(A = QR\), the submission writes

\(\det\left( {A^{T}A} \right) = \det\left( {R^{T}Q^{T}QR} \right) = {\det(R)}^{2},\)

so the determinant is positive.

\section{6.2 Exercise 50}

For the matrix whose \(\left( {i,j} \right)\) entry is \(\min\left( {i,j} \right)\), repeated determinant reduction gives

\(\det(A) = 1.\)

\section{6.3 Exercise 14}

The parallelepiped volume is found from the determinant. Since the displayed vectors are linearly dependent, the result is

\(\text{volume} = 0.\)

\section{7.1 Exercise 12}

For

\(A = \begin{pmatrix}
2 & 0 \\
3 & 4
\end{pmatrix},\)

the characteristic equation gives eigenvalues \(2\) and \(4\). The corresponding eigenvector directions in the work are

\(\lambda = 2:\begin{pmatrix}
{- \frac{2}{3}} \\
1
\end{pmatrix},\quad\lambda = 4:\begin{pmatrix}
0 \\
1
\end{pmatrix}.\)

Thus it diagonalizes to \(D = \begin{pmatrix}
2 & 0 \\
0 & 4
\end{pmatrix}\).

\section{7.1 Exercise 18}

For reflection in a plane, the plane is the \(1\)-eigenspace and has dimension two; its normal direction is the \(- 1\)-eigenspace. With an eigenbasis, the submitted diagonal form is

\(\begin{pmatrix}
1 & 0 & 0 \\
0 & 1 & 0 \\
0 & 0 & {- 1}
\end{pmatrix}.\)

\section{7.1 Exercise 42}

The matrices in \(V\) are written as

\(\begin{pmatrix}
a & b & 0 \\
0 & c & 0 \\
0 & d & e
\end{pmatrix}.\)

The five matrix units in positions \(\left( {1,1} \right),\left( {1,2} \right),\left( {2,2} \right),\left( {3,2} \right),\left( {3,3} \right)\) form the submitted basis, so

\(\dim(V) = 5.\)

\chapter{Part B}

\section{Problem 1}

The proof shows that an alternating bilinear form \(F\) is antisymmetric:

\(F\left( {u + v,u + v} \right) = 0 = F\left( {u,v} \right) + F\left( {v,u} \right).\)

Conversely, antisymmetry gives \(F\left( {u,u} \right) = 0\). Since \(F\left( {e_{1},e_{2}} \right) = 1\), bilinearity then identifies the form with the determinant:

\(F\left( {\begin{pmatrix}
a \\
b
\end{pmatrix},\begin{pmatrix}
c \\
d
\end{pmatrix}} \right) = ad - bc.\)

\section{Problem 2}

Let \(M = \begin{pmatrix}
a & b \\
c & d
\end{pmatrix}\) and let \(T(A) = AM\). The submission verifies linearity. In the ordered basis

\(\mathcal{E} = \left( {E_{11},E_{12},E_{21},E_{22}} \right),\)

it computes

\(\lbrack T\rbrack_{\mathcal{E}} = \begin{pmatrix}
a & c & 0 & 0 \\
b & d & 0 & 0 \\
0 & 0 & a & c \\
0 & 0 & b & d
\end{pmatrix}.\)

The determinant calculation is

\(\det\left( \lbrack T\rbrack_{\mathcal{E}} \right) = \left( {ad - bc} \right)^{2},\)

which agrees with the determinant in any other basis. Finally, for

\(M = \begin{pmatrix}
2 & 1 \\
0 & 2
\end{pmatrix},\)

the only eigenvalue is \(2\), with an eigenspace of dimension \(2 < 4\) for the induced map; therefore the submitted conclusion is that \(T\) is not diagonalizable.

\section{Problem 3}

The construction defines \(z\) by the determinant/cross-product functional in \(\mathbb{R}^{4}\). For the standard choice \(u = e_{1}\), \(v = e_{2}\), \(w = e_{3}\), the work finds

\(z = - e_{4}.\)

It proves that \(z = 0\) exactly when \(u,v,w\) are linearly dependent, that \(z\) is orthogonal to each of \(u,v,w\), and that

\(\det\left( {z,u,v,w} \right) = \left\| z \right\|^{2}.\)

\section{Problem 4}

The response uses the characteristic polynomial to find eigenvalues, and uses similarity to preserve the polynomial. It then applies the displayed eigenvector criterion to decide diagonalizability.

\section{Problem 5}

For a \(2 \times 2\) matrix satisfying \(A^{2} = I\), the submission separates the \(+ 1\) and \(- 1\) eigenvector cases and obtains a diagonal form. The final argument extends this to every dimension by decomposing an arbitrary vector as

\(v = \frac{1}{2}\left( {v + Av} \right) + \frac{1}{2}\left( {v - Av} \right),\)

where the two summands lie in the \(1\)- and \(- 1\)-eigenspaces respectively. Thus \(A\) is diagonalizable.
